\documentclass[12pt,a4paper]{article}

% --- Pacchetti Standard ---
\usepackage[utf8]{inputenc}
\usepackage[italian]{babel}
\usepackage{graphicx}
\usepackage{geometry}
\usepackage{hyperref}
\usepackage{booktabs}
\usepackage{listings} % Per il codice
\usepackage{fancyhdr}
\geometry{
    a4paper,
    margin=2.5cm,
    headheight=15pt, % Risolve il warning di fancyhdr (arrotondato per sicurezza)
    top=3cm          % Aumentiamo un po' il margine superiore per far spazio all'header
}
\pagestyle{fancy}
\fancyhf{} % Pulisce header e footer predefiniti

% --- Configurazione Link ---
\hypersetup{
    colorlinks=true,
    linkcolor=blue,
    urlcolor=cyan,
}

% --- Configurazione Header ---
\fancyhead[L]{F1 Fan Club} % Sinistra
\fancyhead[R]{UniPD} % Destra

% --- Configurazione Footer ---
\fancyfoot[L]{Informatica} % Sinistra
\fancyfoot[R]{Pag. \thepage} % Destra

\renewcommand{\headrulewidth}{0.4pt} % Linea sotto l'header
\renewcommand{\footrulewidth}{0.4pt} % Linea sopra il footer

\begin{document}

% --- Frontespizio ---
\begin{titlepage}
    \centering
    % LOGO UNIPD
    \begin{figure}[h]
        \centering
        \includegraphics[width=0.25\textwidth]{../resources/logo_unipd.png} % Sostituisci con il nome corretto
    \end{figure}
    
    \vspace{0.5cm}
    {\Large \textbf{Università degli Studi di Padova}} \\
    {\large Corso di Laurea in Informatica} \\
    
    \vspace{1.5cm}
    
    % LOGO PROGETTO
    \begin{figure}[h]
        \centering
        \includegraphics[width=0.4\textwidth]{../resources/logo.jpeg}
    \end{figure}
    
    {\Huge \textbf{F1 Fan Club}} \\
    \vspace{0.5cm}
    {\large Relazione per il progetto di Tecnologie Web} \\
    
    \vfill
    
    % --- Tabella Componenti ---
    \begin{table}[h]
        \centering
        \begin{tabular}{lll}
            \textbf{Nome e Cognome} & \textbf{Matricola} & \textbf{Email (studenti.unipd.it)} \\ \midrule
            Giuliano Banchieri     & ------- & giuliano.banchieri@studenti.unipd.it \\
            Alessandro Frison      & 2111932 & alessandro.frison.13@studenti.unipd.it \\
            Riccardo Graziani      & ------- & riccardo.graziani.1@studenti.unipd.it \\
            Chiara Grossele        & ------- & chiara.grossele@studenti.unipd.it \\
            Mattia Oliva Medin     & ------- & mattia.olivamedin@studenti.unipd.it \\
        \end{tabular}
    \end{table}

    \vfill
    {\large Anno Accademico 2025/2026}
\end{titlepage}

\tableofcontents
\newpage

% --- Sezioni Contenuto ---

\section{Introduzione}
Descrizione del progetto "F1 Fan Club". Obiettivi del sito, target di utenza (appassionati di motori) e finalità comunicative.

\newpage

\section{Analisi dei Requisiti}
\subsection{Target di utenza}
Chi sono i visitatori? Tifosi occasionali e appassionati di F1.    
\subsection{SEO}
Strategie per migliorare la visibilità sui motori di ricerca (parole chiave, meta tag, contenuti di qualità).

\section{Progettazione}
Descrizione dell'architettura del sito, delle tecnologie utilizzate e della struttura del database.
\subsection{Schema organizzativo}
Schema organizzativo del sito, inclusi i principali moduli e le interazioni tra essi.
\subsection{Tipi di utente}
Descrizione dei diversi tipi di utenti (es. utente normale, amministratore) e i permessi associati.
\subsection{Funzionalità}
\subsection{Convenzioni interne}
\subsection{Schema del Database}
Descrizione delle tabelle (Utenti, Commenti, Piloti, Classifiche, ...).

\newpage

\section{Realizzazione}
Descrizione delle fasi di sviluppo, delle difficoltà incontrate e delle soluzioni adottate.
\subsection{Struttura e contenuto}
Descrizione delle pagine principali del sito e delle loro funzionalità.
\subsubsection{HTML}
\subsubsection{Popolamento database}
\subsection{Presentazione}
Descrizione dell'aspetto grafico del sito, del layout e delle scelte di design.
\subsubsection{CSS}
\subsubsection{Immagini e icone}
\subsubsection{Font}
\subsubsection{Colori}
\subsection{Comportamento}
Descrizione delle funzionalità interattive del sito.
\subsection{PHP}
\subsubsection{JavaScript}
\subsubsection{Validazione dell'Input}
\subsubsection{Sicurezza}
\subsubsection{Errori di navigazione o del server}
\subsection{Accessibilità e Usabilità}
Misure adottate per garantire l'accessibilità (es. contrasto colori, navigazione da tastiera) e l'usabilità (es. layout intuitivo, tempi di caricamento rapidi).
\subsubsection{Aiuti per lo screen reader}
\subsubsection{Compatibilità}
\subsection{Versionamento del codice}
Uso di Git per il controllo delle versioni e la collaborazione tra i membri del team.

\section{Test e Validazione}
Descrizione dei test effettuati per verificare il corretto funzionamento del sito e la conformità ai requisiti iniziali.
\subsection{Navigabilità ed accessibilità}
\subsection{Falsi positivi}
\subsection{Screen reader}

\newpage

\section{Suddivisione del lavoro}
Breve descrizione dei compiti svolti da ciascun membro del gruppo.

\section{Note}

\end{document}